\worksheettitle
  {Домашняя работа: классы чисел}
  {\modulemark{M03-H} Самостоятельное закрепление}

\studentnameblock

\begin{notebox}[Как выполнять]
  Сначала работайте без основного листа. После выполнения проверьте: деление
  начинается справа, а каждая группа, кроме крайней левой, содержит три цифры.
\end{notebox}

\begin{practicebox}[1. Разделите на классы]
  \begin{enumerate}[label=\textbf{\arabic*.}]
    \item \(3807=\answerblank[42mm].\)
    \item \(92064=\answerblank[42mm].\)
    \item \(705009=\answerblank[46mm].\)
    \item \(6040308=\answerblank[52mm].\)
    \item \(81002040=\answerblank[52mm].\)
    \item \(206000015=\answerblank[56mm].\)
  \end{enumerate}
\end{practicebox}

\begin{practicebox}[2. Разберите число]
  Для числа \(64\,008\,205\) заполните таблицу.
  \begin{center}
  \renewcommand{\arraystretch}{1.4}
  \begin{tabular}{c|ccc}
    \toprule
    класс & миллионы & тысячи & единицы \\
    \midrule
    группа & \answerblank[15mm] & \answerblank[15mm] & \answerblank[15mm] \\
    количество единиц класса & \answerblank[15mm] & \answerblank[15mm] &
      \answerblank[15mm] \\
    \bottomrule
  \end{tabular}
  \end{center}
\end{practicebox}

\begin{practicebox}[3. Объясните]
  Почему группу класса тысяч в числе \(64\,008\,205\) нельзя записать как
  \(8\)?
  \writinglines{2}
\end{practicebox}

\newpage
\begin{practicebox}[4. Соберите числа]
  \begin{enumerate}[label=\textbf{\arabic*.}]
    \item \(28\) миллионов, \(4\) тысячи и \(7\) единиц:
      \answerblank[52mm].
    \item \(6\) миллионов, \(0\) тысяч и \(53\) единицы:
      \answerblank[52mm].
    \item \(104\) миллиона, \(30\) тысяч и \(9\) единиц:
      \answerblank[52mm].
  \end{enumerate}
\end{practicebox}

\begin{warningbox}[5. Исправьте ошибки]
  \begin{enumerate}[label=\textbf{\arabic*.}]
    \item \(8040205=804\,|\,020\,|\,5\).
    \item \(61\) миллион, \(2\) тысячи и \(40\) единиц записали как
      \(61\,2\,40\).
  \end{enumerate}
  Запишите верные варианты и назовите причины ошибок.
  \writinglines{3}
\end{warningbox}

\begin{practicebox}[6. Сравните]
  Поставьте знак \(>\), \(<\) или \(=\). Подчеркните первый слева класс, в
  котором группы различаются.
  \[
    34\,006\,500\ \answerblank[8mm]\ 34\,060\,050,
    \qquad
    7\,205\,008\ \answerblank[8mm]\ 7\,025\,800.
  \]
\end{practicebox}

\begin{checkpointbox}[7. Итог]
  \begin{enumerate}[label=\textbf{\arabic*.}]
    \item \(50703009=\answerblank[52mm].\)
    \item В полученном числе группа класса единиц: \answerblank[16mm],
      количество единиц: \answerblank[16mm].
    \item \(19\) миллионов, \(6\) тысяч и \(4\) единицы:
      \answerblank[48mm].
  \end{enumerate}
  \textbf{Перед сдачей:} проверьте направление деления и длину групп.
\end{checkpointbox}
